\scanpage{274}
\mspage{224}
équivalence, qui permet de décrire $\mathrm{Hot}_W^0(S)$ en termes de $\underline{\mathrm{Prop}}_W S\unsure{\subset|\subseteq}\mathrm{Cat}_{/S}$ et d'un localiseur $W_S^{\unsure{u|v}}$ \uline{stable par} \uuline{tout} \uline{chgt. de base}, permet d'autre part de définir la loi contravariante attendue pour $\unsure{\mathrm{HOT}|\mathrm{Hot}}_W^0(S)$.

D'autre part, quand on échange les rôles de $\deleted{\underline{\mathrm{Prop}}}\underline{\mathrm{Liss}}_W$ et $\mathrm{Prop}_W$, on trouve la théorie complémentaire des $\mathrm{HOT}_W(S)$,
\[
\left\{\begin{aligned}
\mathrm{HOT}_W(S)&\overset{\text{déf}}{=}(\mathrm{Cat}_{/S})(W_S^d)^{-1}
\xleftarrow{\approx}(\underline{\mathrm{Liss}}_W S)(W_S^u)^{-1}\\
\text{sur }\underline{\mathrm{Liss}}_W S,\quad &W_S^d=W_{\unsure{S|J}}^u,
\end{aligned}\right.
\]
la description de $\mathrm{HOT}_W(S)$ en termes de $\underline{\mathrm{Liss}}_W S$ et de $W_S^u$ permet d'en \deleted{établir} \inserted{définir} la variation contravariante en $S$, pour \uline{tte} flèche $f:S'\longrightarrow S$ dans $\mathrm{Cat}$. Tandis que si on veut pouvoir passer aux catégories de \unsure{fractions|foncteurs} sur $\mathrm{Cat}_{/S}$ et $\mathrm{Cat}_{/S'}$ tout entiers \inserted{(via $W_S^d$, $W_{S'}^d$)}, il faut se limiter aux chgts de base $W$-propres.

Du coup, me vient la perplexité si cette complémentarité, voire \inserted{cette} dualité, entre deux ensembles de flèches, les ``lisses'' et les ``propres'', ne serait pas un fait général qui se retrouverait tout au moins dans les catégories de modèles \unsure{proches|parties|pleines} de la \unsure{description} de $\mathrm{HOT}$ ou d'un $\mathrm{Hot}_W$.
% archive sheet number: 274
